\documentclass[../../main.tex]{subfiles}
\begin{document}

\section*{问题}\label{section:ch02-问题}

\begin{exercise}\label{exercise:ch02-2-1}
由
\begin{align*}
f(x)=
\begin{cases}
1, & x\geqslant 0,\\
0, & x<0
\end{cases}
\end{align*}
定义 $f:\mathbb{R}\to\mathbb{R}$。证明：对每个 $x\in\mathbb{R}$，都存在包含 $x$ 的光滑坐标图 $(U,\varphi)$ 和包含 $f(x)$ 的光滑坐标图 $(V,\psi)$，使 $\psi\circ f\circ\varphi^{-1}$ 作为从 $\varphi(U\cap f^{-1}(V))$ 到 $\psi(V)$ 的映射是光滑的，但 $f$ 在我们本章所定义的意义下不光滑。
\end{exercise}

\begin{exercise}\label{exercise:ch02-2-2}
证明命题~\cref{proposition:ch02-乘积映射分量}（映到乘积流形的映射的光滑性）。
\end{exercise}

\begin{exercise}\label{exercise:ch02-2-3}
对下列每个球面之间的映射，计算足够多的坐标表示，以证明它光滑。
\begin{enumerate}[(a)]
\item $\rho_{n}:S^{1}\to S^{1}$ 是 $n\in\mathbb{Z}$ 次幂映射，在复记号下由 $\rho_{n}(z)=z^{n}$ 给出。
\item $\alpha:S^{n}\to S^{n}$ 是对径映射 $\alpha(x)=-x$。
\item $F:S^{3}\to S^{2}$ 由
\begin{align*}
F(w,z)=\bigl(z\overline{w}+w\overline{z},\ \mathrm{i}w\overline{z}-\mathrm{i}z\overline{w},\ |z|^{2}-|w|^{2}\bigr)
\end{align*}
给出，其中我们把 $S^{3}$ 看作 $\mathbb{C}^{2}$ 的子集 $\{(w,z):|w|^{2}+|z|^{2}=1\}$。
\end{enumerate}
\end{exercise}

\begin{exercise}\label{exercise:ch02-2-4}
证明：当 $\overline{B^{n}}$ 被看作一个光滑带边流形时，包含映射 $\overline{B^{n}}\hookrightarrow\mathbb{R}^{n}$ 是光滑的。
\end{exercise}

\begin{exercise}\label{exercise:ch02-2-5}
令 $\mathbb{R}$ 为带其标准光滑结构的实直线，令 $\widetilde{\mathbb{R}}$ 表示同一个拓扑流形，但带例~\cref{example:ch01-另一光滑结构} 中定义的光滑结构。设 $f:\mathbb{R}\to\mathbb{R}$ 是一个在通常意义下光滑的函数。
\begin{enumerate}[(a)]
\item 证明 $f$ 作为从 $\mathbb{R}$ 到 $\widetilde{\mathbb{R}}$ 的映射也是光滑的。
\item 证明 $f$ 作为从 $\widetilde{\mathbb{R}}$ 到 $\mathbb{R}$ 的映射光滑，当且仅当每当 $n$ 不是 $3$ 的整数倍时 $f^{(n)}(0)=0$。
\end{enumerate}
\end{exercise}

\begin{exercise}\label{exercise:ch02-2-6}
设 $P:\mathbb{R}^{n+1}\setminus\{0\}\to\mathbb{R}^{k+1}\setminus\{0\}$ 是一个光滑函数，并假设对某个 $d\in\mathbb{Z}$，对所有 $\lambda\in\mathbb{R}\setminus\{0\}$ 和 $x\in\mathbb{R}^{n+1}\setminus\{0\}$ 有 $P(\lambda x)=\lambda^{d}P(x)$。（这样的函数称为 $d$ 次齐次的。）证明由 $\widetilde{P}([x])=[P(x)]$ 定义的映射 $\widetilde{P}:\mathbb{RP}^{n}\to\mathbb{RP}^{k}$ 良定义且光滑。
\end{exercise}

\begin{exercise}\label{exercise:ch02-2-7}
设 $M$ 是一个非空的、维数 $n\geqslant 1$ 的带边或无边光滑流形。证明向量空间 $C^{\infty}(M)$ 是无限维的。[提示：证明若 $f_{1},\dots,f_{k}$ 是 $C^{\infty}(M)$ 中支集非空且两两不相交的元素，则它们线性无关。]
\end{exercise}

\begin{exercise}\label{exercise:ch02-2-8}
由
\begin{align*}
F\bigl(x^{1},\dots,x^{n}\bigr)=\bigl[x^{1},\dots,x^{n},1\bigr]
\end{align*}
定义 $F:\mathbb{R}^{n}\to\mathbb{RP}^{n}$。证明 $F$ 是一个到 $\mathbb{RP}^{n}$ 的一个稠密开子集的微分同胚。对由
\begin{align*}
G\bigl(z^{1},\dots,z^{n}\bigr)=\bigl[z^{1},\dots,z^{n},1\bigr]
\end{align*}
定义的 $G:\mathbb{C}^{n}\to\mathbb{CP}^{n}$ 也作同样证明（见习题~\cref{exercise:ch01-1-9}）。
\end{exercise}

\begin{exercise}\label{exercise:ch02-2-9}
给定一个一元、复系数、不恒为零的多项式 $p$，证明存在唯一的光滑映射 $\widetilde{p}:\mathbb{CP}^{1}\to\mathbb{CP}^{1}$ 使下图交换，其中 $\mathbb{CP}^{1}$ 是 $1$ 维复射影空间，$G:\mathbb{C}\to\mathbb{CP}^{1}$ 是习题~\cref{exercise:ch02-2-8} 中的映射：
即 $G\circ p=\widetilde{p}\circ G$。
（用于第~465 页。）
\end{exercise}

\begin{exercise}\label{exercise:ch02-2-10}
对任意拓扑空间 $M$，令 $C(M)$ 表示连续函数 $f:M\to\mathbb{R}$ 所成代数。给定连续映射 $F:M\to N$，由 $F^{*}(f)=f\circ F$ 定义 $F^{*}:C(N)\to C(M)$。
\begin{enumerate}[(a)]
\item 证明 $F^{*}$ 是一个线性映射。
\item 假设 $M,N$ 是光滑流形。证明 $F:M\to N$ 光滑，当且仅当 $F^{*}(C^{\infty}(N))\subset C^{\infty}(M)$。
\item 假设 $F:M\to N$ 是光滑流形之间的一个同胚。证明它是微分同胚，当且仅当 $F^{*}$ 限制为一个从 $C^{\infty}(N)$ 到 $C^{\infty}(M)$ 的同构。
\end{enumerate}
[注：这个结果表明，在某种意义上，$M$ 的整个光滑结构都编码在子集 $C^{\infty}(M)\subset C(M)$ 中。事实上，有些作者把拓扑流形 $M$ 上的光滑结构定义为 $C(M)$ 的一个具有某些性质的子代数；例如见 [Nes03]。]（用于第~75 页。）
\end{exercise}

\begin{exercise}\label{exercise:ch02-2-11}
假设 $V$ 是一个维数 $n\geqslant 1$ 的实向量空间。定义 $V$ 的射影化，记作 $P(V)$，为 $V$ 的 $1$ 维线性子空间所成的集合，带由映射 $\pi:V\setminus\{0\}\to P(V)$（把 $x$ 映到它张成的子空间）诱导的商拓扑。（于是 $P(\mathbb{R}^{n})=\mathbb{RP}^{n-1}$。）证明 $P(V)$ 是一个拓扑 $(n-1)$ 流形，并且有唯一的光滑结构，使对 $V$ 的每个基 $(E_{1},\dots,E_{n})$，由
\begin{align*}
E\bigl([v^{1},\dots,v^{n}]\bigr)=\bigl[v^{i}E_{i}\bigr]
\end{align*}
定义的映射 $E:\mathbb{RP}^{n-1}\to P(V)$（其中方括号表示等价类）是一个微分同胚。（用于第~561 页。）
\end{exercise}

\begin{exercise}\label{exercise:ch02-2-12}
叙述并证明习题~\cref{exercise:ch02-2-11} 对复向量空间的类比。
\end{exercise}

\begin{exercise}\label{exercise:ch02-2-13}
假设 $M$ 是一个拓扑空间，具有性质：对 $M$ 的每个带指标开覆盖 $\mathcal{X}$，都存在一个从属于 $\mathcal{X}$ 的单位分解。证明 $M$ 是仿紧的。
\end{exercise}

\begin{exercise}\label{exercise:ch02-2-14}
假设 $A,B$ 是光滑流形 $M$ 的两个不相交闭子集。证明存在 $f\in C^{\infty}(M)$，使对所有 $x\in M$ 有 $0\leqslant f(x)\leqslant 1$，$f^{-1}(0)=A$，且 $f^{-1}(1)=B$。
\end{exercise}

\end{document}